\subsection{Falsifiability}
  \label{subsec:falsifiability}

  The effective saturation framework makes a number of concrete,
  observationally testable predictions.
  Because the same universal scale $a_\star$ governs both
  galactic dynamics and late-time kinematic inference,
  the model can be falsified at multiple, independent levels.

  \paragraph{(i) Absence of environment dependence in $H_0$.}

    The framework predicts a correlation between the locally inferred
    Hubble parameter and the surrounding large-scale density contrast,
    as expressed in Eq.~\eqref{eq:hloc}.
    Environment-stratified analyses of distance-ladder measurements
    should reveal a statistically significant trend between inferred $H_0$
    and void proximity or density contrast.
    If no such correlation is detected within observational precision,
    or if the inferred $H_0$ proves statistically independent of environment,
    the proposed mechanism would be strongly disfavored.

  \paragraph{(ii) Redshift scaling of the offset.}

    The predicted deviation $\Delta H/H$ is expected to decrease
    with increasing redshift due to environmental averaging
    along extended lines of sight.
    A persistent offset of comparable magnitude at higher redshift,
    inconsistent with large-scale structure dilution,
    would contradict the core mechanism.
    Conversely, a measured redshift scaling incompatible with
    the functional dependence implied by Eq.~\eqref{eq:hloc}
    would also falsify the framework.

  \paragraph{(iii) Inconsistency of the universal scale.}

    The acceleration scale $a_\star$ is fixed from galaxy rotation curves.
    Cosmological applications do not introduce a new free parameter.
    If independent cosmological analyses required a significantly
    different effective scale to reproduce any environment-dependent
    expansion signal, the internal consistency of the framework
    would be violated.
    The model therefore predicts cross-consistency between galactic
    and cosmological determinations of the same parameter.

  \paragraph{(iv) Breakdown of galactic saturation behavior.}

    The model predicts systematic low-acceleration behavior tied to
    a single universal scale across high- and low-surface-density
    galaxies.
    A statistically significant population of extremely diffuse
    systems that cannot be reconciled with a single $a_\star$,
    or that require environment-dependent tuning of the scale,
    would invalidate the bounded-response interpretation.

  \paragraph{(v) Standard-model sufficiency.}

    Finally, if improved $\Lambda$CDM-based modeling of feedback,
    selection effects, and local structure fully resolves both
    galaxy-scale phenomenology and the Hubble tension without
    invoking any environment-dependent scaling tied to a universal
    acceleration threshold, the explanatory necessity of the present
    mechanism would disappear.

    These criteria render the framework empirically vulnerable.
    It does not merely accommodate the observed anomalies,
    but constrains their amplitude, scaling, redshift dependence,
    and cross-scale consistency.
    Failure of any of the above conditions would require
    revision or rejection of the proposed mechanism.
